%\needspace{4\baselineskip}
\color{uniblue}{\subsection{Устройство резервирования при отказе выключателя}}
\color{black}

\begin{enumerate}[label=\arabic{section}.\arabic{subsection}.\arabic{enumi}, labelsep=4pt, leftmargin=0pt, itemindent=57pt, itemsep=0pt, parsep=0pt] % если что можно крутить itemsep


\item
В устройстве предусмотрена функция резервирования при отказе выключателя (УРОВ). Данная 
функция предназначена для ликвидации повреждения в случае отказа выключателя поврежденного 
присоединения.

\item
Логика работы УРОВ выполнена с контролем протекания тока через выключатель. Значение по 
току срабатывания задается параметром \textbf{«Iср УРОВ В1(ЛВ)ВН»}. Функция УРОВ действует с 
регулируемой задержкой по времени \textbf{«Tср повт. УРОВ»} на повторное отключение своего выключателя, 
с задержкой \textbf{«Tср УРОВ»} на отключение смежных выключателей. 

\item
Предусмотрен подхват по току, включение которого производится программной накладкой 
\textbf{«УРОВ подхв.»}. 

\item
Действие УРОВ на смежные выключатели может быть выполнено с контролем сигнала РПВ. 
Ввод контроля сигнала РПВ в логику работы УРОВ осуществляется программной накладкой \textbf{«УРОВ 
контр. выкл.»}. 


\begin{figure}[H]
\centering
\includegraphics[width=1\textwidth,height=1\textheight,keepaspectratio]{\fbpath/УРОВ/img/UROV.pdf}
\captionof{figure}{Функциональный блок УРОВ}\label{pict:UROV}
\end{figure}




\small
\setlength{\extrarowheight}{0.05cm} % высота ячеек в таблице
\begin{longtable}
{|>{\centering\arraybackslash}m{4.5cm}|>{\centering\arraybackslash}m{4.5cm}|>{\centering\arraybackslash}m{3cm}|>{\centering\arraybackslash}m{3.5cm}|}
\caption{Перечень уставочных параметров УРОВ\hfill\vspace{-0.5\baselineskip}}\label{ksn:tbl2} \\ 
\hhline{|-|-|-|-|}
\rowcolor{gray!30}
Наименование уставки & Диапазон & Шаг &  Значение по умолчанию \\ 
\hhline{|-|-|-|-|}
\endfirsthead

\caption*{\hspace{3pt}\emph{Продолжение таблицы \ref{ksn:tbl2}\hfill\vspace{-0.5\baselineskip}}} \\ % сделано по ГОСТ 2.105 п.6.8.7
\hhline{|-|-|-|-|}
\rowcolor{gray!30}
Наименование уставки & Диапазон & Шаг &  Значение по умолчанию \\  
\hhline{|-|-|-|-|}
\endhead
\hhline{|-|-|-|-|}

\endfoot
\endlastfoot

%%%%%%%%%%%%содержимое%%%%%%%%%%%%%%

\end{longtable}
\setlength{\extrarowheight}{0.0cm} % возврат параметра высота ячеек в таблице
\normalsize























\end{enumerate} 